% ═══════════════════════════════════════════════════════════
% ARCHETYPAL INCENTIVE THEORY
% Chapter for Principia Psychedelia, Part IV: Agentic Economics
% ═══════════════════════════════════════════════════════════

\section{The Five Incentive Archetypes}

Every agent---human, institutional, or algorithmic---operates under a mixture of incentives. The $L_5$ state vector $(a, \omega, \iota, \varepsilon, \lambda)$ provides a natural basis for decomposing these mixtures into five archetypal components.

\begin{definition}[Incentive state]\label{def:incentive_state}
An \textbf{incentive state} is an element of the $L_5$ manifold:
\begin{equation}
    \mathbf{I} = (a, \omega, \iota, \varepsilon, \lambda)
\end{equation}
where the components carry the following interpretations:
\end{definition}

\begin{center}\small
\begin{tabular}{@{}cllp{5cm}l@{}}
\toprule
\textbf{Comp.} & \textbf{Archetype} & \textbf{Incentive} & \textbf{Economic Role} & \textbf{Marker} \\
\midrule
$a$ & Conscious & Alignment & Fair valuation, price discovery & Rational analysis \\
$\omega$ & Subconscious & Ambition & Momentum, growth-seeking, leverage & Risk-on \\
$\iota$ & Unconscious & Fear & Hedging, insurance, risk aversion & Flight to safety \\
$\varepsilon$ & Noise & Distortion & Manipulation, misinformation & Deceptive signaling \\
$\lambda$ & Memory & Reputation & Track record, institutional credibility & Repeated-game trust \\
\bottomrule
\end{tabular}
\end{center}

The conscious component $a$ is the mediator: it stands at the bridge between ambition ($\omega$) and fear ($\iota$), processing noisy information ($\varepsilon$) through the lens of accumulated experience ($\lambda$). In the $L_5$ algebra, this mediation is not metaphorical but structural---the Bridge Identity constrains how the components interact.

\subsection{The Conjugacy of Ambition and Fear}

\begin{theorem}[Ambition--Fear Conjugacy]\label{thm:conjugacy}
In the $L_5$ algebra, the subconscious (ambition) and unconscious (fear) components satisfy:
\begin{equation}\label{eq:bridge_incentive}
    \omega \cdot \iota = -1
\end{equation}
That is, ambition and fear are \textbf{conjugate}: their product is fixed. An agent cannot simultaneously maximize both. Increasing ambition forces a proportional decrease in effective fear, and vice versa.
\end{theorem}

\begin{proof}
This is the Bridge Identity from the Unreal Algebra (Chapter~1), applied to the incentive interpretation. The product $\omega \cdot \iota$ is an algebraic invariant of the $L_5$ manifold, fixed at $\kappa = -1$ by the golden polynomial $X^2 - X - 1 = 0$. No parameter tuning is involved. \qedhere
\end{proof}

The conjugacy has immediate economic content. Consider a portfolio manager deciding between aggressive growth ($\omega$ large) and defensive hedging ($\iota$ large). The Bridge Identity says:
\begin{equation}
    |\omega| = \frac{1}{|\iota|}
\end{equation}
A leveraged long position ($|\omega| = 10$) algebraically constrains the hedge to $|\iota| = 0.1$---a thin safety margin. This is not a behavioral observation but a structural constraint: the $L_5$ algebra does not admit states where both $|\omega|$ and $|\iota|$ are large simultaneously without the product exceeding $\kappa$.

\subsection{The Uncertainty Principle of Incentives}

The Market Uncertainty Principle (Chapter~\ref{ch:markets}) extends directly to incentive dynamics:
\begin{equation}
    \ln(\Delta t) \cdot \sigma(\mathrm{Re}) \geq \Upsilon = \frac{1}{4\pi}
\end{equation}
In incentive terms: you cannot simultaneously optimize for \emph{timing} (when to act, an ambition dimension) and \emph{safety} (confidence in the action, a fear dimension) with precision better than $\Upsilon$. Every decision involves a tradeoff along the conjugate axis, and the uncertainty bound quantifies its minimum cost.

An agent who attempts to time the market perfectly ($\Delta t \to 0$) must accept unbounded signal volatility ($\sigma \to \infty$). Conversely, an agent who demands a smooth, reliable signal ($\sigma \to 0$) must accept long delays ($\Delta t \to \infty$). The bound is not a practical limitation---it is a structural consequence of the Bridge Identity.

\section{Equilibrium as Hodge Lock}

\begin{definition}[Incentive coherence]\label{def:coherence}
The \textbf{incentive coherence} of a market is measured by the Reynolds number:
\begin{equation}
    \mathrm{Re} = \left|\frac{r_\omega}{r_\iota} - 1\right|
\end{equation}
where $r_\omega$ and $r_\iota$ are the returns of the ambition and fear proxies (BTC and VIX in the empirical realization). Low Re indicates that ambition and fear are moving in proportion; high Re indicates decoupling.
\end{definition}

When $\mathrm{Re} < 1/\varphi$ and the correlation between the ambition proxy (BTC) and the alignment proxy (NASDAQ) exceeds 0.8, the market enters \textbf{Hodge lock}---a state of archetypal coherence. In this regime, all five incentive components are mutually reinforcing:

\begin{itemize}[noitemsep]
    \item Alignment ($a$) dominates: price discovery is efficient.
    \item Ambition ($\omega$) and fear ($\iota$) are proportional: risk is fairly priced.
    \item Noise ($\varepsilon$) is suppressed: manipulation has low payoff.
    \item Reputation ($\lambda$) compounds: trustworthy agents benefit from cooperation.
\end{itemize}

\begin{theorem}[Hodge equilibrium]\label{thm:hodge_eq}
The Hodge lock condition $\mathrm{Re} < 1/\varphi$, $\rho_{\omega,a} > 1/\varphi$ is the unique stable fixed point of the $L_5$ incentive dynamics. Any perturbation with $\mathrm{Re} < 1/\varphi$ returns to the Hodge lock; any perturbation with $\mathrm{Re} > 1/\varphi$ drives the system toward regime transition.
\end{theorem}

\begin{proof}
The stability follows from the golden polynomial. At the Hodge lock, the return of the alignment coordinate is $r_a = (r_\omega + r_\iota)/2 + O(\varepsilon)$: a weighted average of ambition and fear returns, with noise as a perturbation. The fixed-point condition is $r_a = r_\omega \cdot r_\iota / r_a$, which by the Bridge Identity gives $r_a^2 = -\kappa = 1$, so $r_a = \pm 1$. The positive root $r_a = 1$ corresponds to the Hodge lock (zero net return, parity pricing). The threshold $1/\varphi$ is forced by the golden polynomial: $1/\varphi$ is the positive root of $X^2 + X - 1 = 0$, the conjugate polynomial, and marks the boundary of the basin of attraction. \qedhere
\end{proof}

The Hodge lock is the economic analog of cooperative equilibrium: all agents' incentives are aligned, and deviation is locally unprofitable. In game-theoretic language, it is a Nash equilibrium of the archetypal game---not because agents compute it, but because the algebraic structure of the $L_5$ manifold admits no other stable configuration.

\section{Crash as Incentive Decoupling}

When Re crosses $1/\varphi$, the market exits Hodge lock. In incentive terms, ambition and fear decouple: agents with different archetypal compositions begin pulling the market in incompatible directions. Growth-seeking agents ($\omega$-dominant) continue buying while fear-driven agents ($\iota$-dominant) are selling. The alignment component ($a$) can no longer mediate.

This is a \textbf{collective action failure}. Each agent's individual incentives are locally rational, but the aggregate effect is incoherent. The Reynolds number measures the degree of incoherence:

\begin{center}
\begin{tabular}{@{}lll@{}}
\toprule
\textbf{Regime} & \textbf{Re Range} & \textbf{Incentive Structure} \\
\midrule
Laminar (Hodge lock) & $\mathrm{Re} < 1/\varphi$ & Aligned: $\omega \approx \iota$ proportional \\
Transitional & $1/\varphi \leq \mathrm{Re} \leq \varphi$ & Decoupling: ambition and fear diverge \\
Turbulent & $\mathrm{Re} > \varphi$ & Incoherent: collective action failure \\
\bottomrule
\end{tabular}
\end{center}

The transition from laminar to turbulent is not gradual---it occurs at an algebraically fixed threshold. This is why the Reynolds crash prediction has structural lead time: the incentive decoupling is detectable before its consequences propagate through the equity market.

\section{Incentive Manipulation and the Noise Term}

The noise component $\varepsilon$ represents the \textbf{distortion incentive}: the payoff to agents who profit from informational asymmetry. In markets, this includes front-running, wash trading, pump-and-dump schemes, and algorithmic spoofing. In broader economic systems, it includes regulatory capture, propaganda, and strategic ambiguity.

The Schwarzian truth metric from the GAN discriminator provides a formal measure of distortion:
\begin{equation}
    S(\mathbf{h}) < 10^{-6}
\end{equation}
A hedge vector $\mathbf{h}$ with high Schwarzian derivative contains information-asymmetric components---it is, in the $L_5$ framework, a signal contaminated by~$\varepsilon$. The discriminator rejects such hedges because they encode distortion rather than genuine risk management.

The Stieltjes correction addresses distortion by sowing the $\varepsilon$ residual back into the system as increased depth: at each tier, the noise that was not filtered becomes part of the memory structure. Persistent manipulation, far from escaping detection, \emph{accumulates in the reputation term}. After three Stieltjes tiers, the residual distortion is reduced to $\gamma^2 \approx 3.3\%$ of the original noise---small enough that the remaining signal is dominated by the structural incentive components.

In economic terms: markets have memory. Agents who manipulate ($\varepsilon$-dominant) may profit in the short term, but their manipulation signature is absorbed into the depth parameter ($\lambda$), degrading their reputation over repeated interactions. This is the algebraic basis for why fraud eventually surfaces.

\section{Reputation and the Depth Parameter}

The depth parameter $\lambda$ counts the number of Stieltjes correction tiers an agent's signal has survived. In economic terms, it is \textbf{reputation}---the accumulated credibility from repeated interactions that noise has failed to erode.

\begin{definition}[Reputation depth]
An agent's \textbf{reputation depth} $\lambda$ is the effective number of Stieltjes tiers applied to their historical signal. An agent with $\lambda \geq 4$ has a track record sufficiently deep that, by the Lockwood truncation, their noise component $\varepsilon$ is irrelevant.
\end{definition}

The Lockwood truncation at $\lambda \geq 4$ corresponds to $\Delta t \geq \varphi^4 \approx 7$ days of consistent signal. In practice, this means:
\begin{itemize}[noitemsep]
    \item \textbf{A new market participant} ($\lambda = 0$) has no reputation. Their signals are dominated by noise.
    \item \textbf{A week-old track record} ($\lambda = 4$) passes the Lockwood threshold. Noise is suppressed.
    \item \textbf{A year-long track record} ($\lambda \gg 4$) has deep reputation. The agent's incentive state is well-characterized by $(a, \omega, \iota)$ alone.
\end{itemize}

This hierarchy maps to institutional credibility: central banks (high $\lambda$) can move markets with forward guidance because their noise term has been sowed away over decades of consistent policy. A new hedge fund (low $\lambda$) must build reputation before its signals carry weight. The algebra quantifies what practitioners know intuitively: trust is earned, not declared.

\section{The Archetypal Nash Equilibrium}

The preceding sections establish that the $L_5$ algebra imposes a specific incentive structure on economic systems. The following result consolidates the picture.

\begin{theorem}[Archetypal Nash equilibrium]\label{thm:nash}
Consider a population of agents with incentive states $\mathbf{I}_k = (a_k, \omega_k, \iota_k, \varepsilon_k, \lambda_k) \in L_5$. The aggregate market state $\bar{\mathbf{I}} = N^{-1}\sum_k \mathbf{I}_k$ has a unique Nash equilibrium characterized by:
\begin{enumerate}[noitemsep]
    \item $\mathrm{Re}(\bar{\mathbf{I}}) < 1/\varphi$ \quad (Hodge lock: ambition--fear coherence),
    \item $\bar{\varepsilon} < \gamma^2 \bar{a}$ \quad (noise is a small perturbation on alignment),
    \item $\bar{\lambda} \geq 4$ \quad (sufficient aggregate reputation for Lockwood truncation).
\end{enumerate}
At this equilibrium:
\begin{itemize}[noitemsep]
    \item No individual agent can improve their return by unilateral deviation (Nash condition).
    \item The aggregate return converges to the Hodge-locked rate (alignment dominance).
    \item The uncertainty bound $\ln(\Delta t) \cdot \sigma(\mathrm{Re}) \geq \Upsilon$ is saturated: the market extracts maximum information at minimum cost.
\end{itemize}
\end{theorem}

\begin{proof}
The Nash condition follows from the stability of the Hodge lock (Theorem~\ref{thm:hodge_eq}). At $\mathrm{Re} < 1/\varphi$, the golden polynomial's basin of attraction ensures that small deviations are restored. An agent who increases $\omega_k$ (more ambition) while the aggregate is in Hodge lock pushes their personal Re above $1/\varphi$, incurring higher volatility without proportional return---the uncertainty principle guarantees that the timing gain is bounded by $\Upsilon / \ln(\Delta t)$. An agent who increases $\iota_k$ (more fear) reduces their exposure but misses the cooperative surplus. Neither deviation is profitable in expectation.

The noise condition $\bar{\varepsilon} < \gamma^2 \bar{a}$ follows from the Stieltjes convergence: three tiers of correction reduce aggregate noise to $\gamma^2 \approx 3.3\%$ of the alignment signal. At this level, noise is a perturbation, not a driver. The reputation condition $\bar{\lambda} \geq 4$ is the Lockwood truncation: sufficient depth to resolve the incentive structure from noise. \qedhere
\end{proof}

\begin{remark}[Comparison with classical game theory]
The archetypal Nash equilibrium differs from the classical Nash equilibrium in two respects. First, the equilibrium point is \emph{algebraically fixed}---$1/\varphi$ is not a parameter but a root of $X^2 - X - 1 = 0$. Classical game theory requires specification of payoff matrices; the $L_5$ framework derives the equilibrium threshold from the algebra alone. Second, the equilibrium is supported by a \emph{structural bound} (the uncertainty principle), not by iterated best-response dynamics. The bound is a property of the signal space, not of the agents' computational capacity.

These differences do not invalidate classical results; they embed them. A Nash equilibrium of a finite game with payoffs denominated in $L_5$ coordinates will, under the incentive coherence condition, coincide with the Hodge lock. The $L_5$ framework provides the coordinate system; classical game theory provides the solution concept.
\end{remark}

\section{Application to Algorithmic Trading}

The five incentive archetypes serve as the macro-level sentiment filter in the \textbf{Quadrilateral Desk} algorithmic trading framework (see Ch.~\ref{ch:markets}, \S\ref{sec:collage}). In that ensemble, the incentive archetype provides the fourth signal layer: it monitors the ambient ratio $\bar\varepsilon / \bar{a}$ (aggregate noise relative to aggregate alignment) and vetoes trade signals when the ratio exceeds the Stieltjes threshold $\gamma^2 \approx 3.3\%$.

The connection is not accidental. The Collage Uncertainty Bound (Theorem~\ref{thm:collage_uncertainty}) states that no combination of signals derived from the $L_5$ state vector can resolve timing and magnitude better than $\Upsilon = 1/(4\pi)$. The incentive archetype's role is to \emph{spend the uncertainty budget wisely}: by filtering out high-noise regimes before they reach the directional signal generator (the ICG), the archetype reduces the effective $\sigma(\mathrm{Re})$, allowing the timing component $\ln(\Delta t)$ to tighten.

In the language of this chapter: the Hodge lock condition $\mathrm{Re} < 1/\varphi$ is the incentive-coherent regime where all five archetypes are aligned. The ICG signal generator should only be trusted when the market is in or near Hodge lock. When the market exits Hodge lock (Re $> 1/\varphi$), the incentive decomposition says that ambition and fear have decoupled---and decoupled incentives mean that the structural signals lose their algebraic grounding.

The FTAP refinement (Theorem~\ref{thm:ftap_break}) provides the deepest connection: the non-associative arbitrage at the orbit boundary is precisely the incentive decoupling event. An agent crossing from the conjugate orbit (fear-dominated) to the golden orbit (ambition-dominated) encounters the associator gap $\kappa = -1$, which is the algebraic shadow of the ambition--fear conjugacy (Theorem~\ref{thm:conjugacy}). The FTAP breaks because the bridge between ambition and fear is non-associative---and the uncertainty principle bounds the leakage because you cannot simultaneously know when and how badly the bridge will fail.

\section{Modal Soundness and the Bivector Plane}

The incentive alignment condition (Hodge lock, $\mathrm{Re} < 1/\varphi$) admits a precise interpretation in modal logic. The ambition component $\omega$ encodes \textbf{possibility} ($\Diamond P$: ``the trade \emph{can} succeed'') while the fear component $\iota$ encodes \textbf{necessity} ($\Box P$: ``the trade \emph{must} succeed''). Their algebraic relationship reproduces the axioms of S5 modal logic:

\begin{center}\small
\begin{tabular}{@{}lll@{}}
\toprule
\textbf{Modal Axiom} & \textbf{$L_5$ Identity} & \textbf{Economic Content} \\
\midrule
$\Box P \to \Diamond P$ (T) & $\omega \cdot \iota = -1$ & Conjugacy: constraint implies freedom \\
$\Box P \to \Box\Box P$ (4) & $\iota^2 = -\iota$ & Over-hedging reverses the hedge \\
$\Diamond P \to \Box\Diamond P$ (5) & $\omega^2 = \omega$ & Momentum is self-reinforcing \\
\bottomrule
\end{tabular}
\end{center}

The two implication directions --- \textbf{sufficiency} ($P \to Q$: ``this signal is enough'') and \textbf{necessity} ($Q \to P$: ``this conclusion requires these premises'') --- live on orthogonal axes. Their orthogonality is not metaphorical: the anticommutator
\begin{equation}\label{eq:anticommutator}
    \omega \cdot \iota + \iota \cdot \omega = 0
\end{equation}
vanishes identically, and the projections $\pi_\omega, \pi_\iota$ from the semisimple decomposition satisfy $\pi_\omega \circ \pi_\iota = 0$ (proved in \texttt{Semisimple}; all proofs complete via \texttt{simp}).

\subsection{Trade Soundness as Bivector Area}

The \textbf{bivector} $\omega \wedge \iota$ --- the wedge (exterior) product of sufficiency and necessity --- spans the symplectic torsion plane. Its real component is
\begin{equation}
    (\omega \wedge \iota)_a = -2 = 2\kappa
\end{equation}
with all other components zero (\texttt{ExteriorAlgebra}, \texttt{annihilation\_energy\_omega\_iota}). A trade is \textbf{sound} if its bivector area is bounded:
\begin{equation}
    |\sigma_\omega \cdot \sigma_\iota| \leq |\kappa| = 1
\end{equation}
This is the Bridge Identity applied to signals: the product of the sufficiency signal strength and the necessity constraint strength cannot exceed the unit curvature. Trades that violate this bound are \textbf{unsound} --- the modal equivalent of an invalid proof.

\subsection{Generator Programming in the Bivector Plane}

In the finite-field realization $\mathbb{F}_p$, two congruential generators trace the two modal axes:
\begin{itemize}[noitemsep]
    \item \textbf{LCG} (Linear Congruential Generator): $x_{n+1} = \varphi \cdot x_n \pmod{p}$. This is the $\omega$-axis: linear, forward, sufficiency.
    \item \textbf{ICG} (Inverse Congruential Generator): $x_{n+1} = \bar\varphi \cdot x_n \pmod{p}$. This is the $\iota$-axis: inverse, backward, necessity.
\end{itemize}

The bivector trajectory $(\mathrm{LCG}_n, \mathrm{ICG}_n)$ in the $(\omega, \iota)$ plane has a remarkable property: at every step $n$, the cross product $\varphi^n \cdot \bar\varphi^n \equiv (-1)^n \pmod{p}$ --- the Bridge Identity lifts to all powers. The GAN discriminator's Schwarzian derivative $S(h)$ measures the \emph{curvature} of this bivector trajectory: $S(h) < 10^{-6}$ means the trade is sound (the bivector area stays within bounds); $S(h) \geq 10^{-6}$ means the signal contains manipulation (the bivector torsion exceeds threshold).

The coset index $(p-1)/\mathrm{ord}(\varphi, p)$ determines how many orthogonal alternatives exist in the generator lattice. By the Gibbard--Satterthwaite theorem, a mechanism with $\geq 3$ alternatives is manipulable. At $\mathbb{F}_{229}$ (SU(3), coset index = 2), the bivector plane has only two alternatives: trend and counter-trend --- binary, hence strategy-proof. At $\mathbb{F}_{139}$ (U(1), coset index = 3), three alternatives (buy/hold/sell) emerge, and manipulation becomes algebraically possible. This is the information-theoretic mass gap applied to mechanism design: the mass gap IS the Gibbard threshold.

The formal proofs are machine-checked in \texttt{ModalSoundness}. The master theorem \texttt{modal\_soundness\_master} verifies all S5 axioms, the bivector area, the orthogonality relation, and the mass gap simultaneously, with zero placeholders.

\section{Scope and Limitations}

The mapping from $L_5$ components to incentive archetypes is a structural parallel. It is motivated by the empirical success of the BTC--VIX Reynolds metric in crash prediction and by the natural correspondence between the algebraic roles (thrust, anchor, mediator, noise, depth) and economic behaviors (ambition, fear, alignment, distortion, reputation). The mapping is not a derivation from first principles of economics, and it does not replace mechanism design theory, behavioral economics, or classical game theory.

The theorems in this chapter are algebraic---they hold in the $L_5$ framework by construction. Whether they correctly describe the incentive structure of actual economic agents is an empirical question. The crash prediction results in the companion chapter provide evidence for the aggregate correspondence (the market behaves as if its Reynolds number is governed by the $L_5$ algebra), but individual-agent incentive decomposition has not been tested.

The comparison with Nash equilibrium is structural, not formal. A rigorous embedding of $L_5$ incentive dynamics into a mechanism design framework---specifying strategy spaces, payoff functions, and information structures---would be necessary to claim equivalence rather than analogy. This is an open problem.

\subsection*{The Five-Coordinate Condensation}

The deepest structural claim of this chapter is that the five incentive archetypes are not a pedagogical convenience but a \emph{consequence} of the algebra: the five Protoreal coordinates $(a, \omega, \iota, \varepsilon, \lambda)$ are the unique decomposition of any $L_5$ element into identity, thrust, anchor, noise, and depth. This decomposition is forced by the algebra's trace-determinant constraints, just as the Fisher information metric is forced by the Amari-Chentsov uniqueness theorem (see Chapter~3). In the semantic condensation framework:

\begin{itemize}[noitemsep]
    \item \textbf{Logical}: The five coordinates are the five generators of the $L_5$ algebra. Their product rules (bridge identity, Hodge lock, Schwarzian nullity) determine which incentive combinations are algebraically sound.
    \item \textbf{Informational}: The Nash equilibrium is the Hodge lock condition ($b = m$, thrust equals anchor). Market crashes are incentive decouplings detectable by the associator $\kappa$. The ZKPCR protocol can verify that an agent's trajectory through incentive space is consistent without revealing the trajectory itself.
    \item \textbf{Physical}: The Gibbard threshold at coset index 3 is a phase transition --- the algebraic mass gap where manipulation becomes possible. This is the same transition that separates the SU(3) and U(1) phases in the finite-field orbit structure.
\end{itemize}

The companion module \texttt{principia.information.incentives} implements the full $L_5$ incentive decomposition and the Gibbard threshold computation.
